\scanpage{312}\mspage{36}
que si $\mathcal F$ est une de ces \deleted{catégories} \inserted{classes}, la
condition
\[
\mathcal F\subset\deleted{\unsure{\mathcal F}}\mathrm{Fib}_W
\]
équivaut à celle-ci :

Pour \deleted{\unsure{tout}} diagramme cartésien
\diagram{\begin{tikzcd}[ampersand replacement=\&,column sep=normal,row sep=small]
\& X\arrow[d,dash,"f"'] \& X'\arrow[l,"q"']\arrow[d,"f'"]\\
(4.2.2) \& S \& S'\arrow[l,"p"]\text{,}
\end{tikzcd}}
avec $f\in\mathcal F$ et $p\in W$, on \unsure{á|a} $q\in W$. (NB, on
a aussi $f'\in\mathcal F$, par \inserted{la} stabilité rappelée plus haut.)

4.\deleted{\unsure{2}}3. Voici donc \inserted{(cf plus bas)} l'axiome pertinent, $\mathrm{loc}(7)$
\deleted{\unsure{avec}} en trois aspects $a,b,c$\textsuperscript{*}. On
% slanted
\margin{j'y ai rajouté finalement une quatrième, $\deleted{\unsure{\mathrm F}},\unsure{\varphi|q}$ (comme ``\unsure{fibration|fibrations}'')}
a vu (et on verra) que l'on a les
implications
\[
a\Longrightarrow b\Longrightarrow c.\,[\Longrightarrow\unsure{\mathrm F|\varphi}].
\]
Quand on référera à $\mathrm{Loc}(7)$ sans autre précision,
il sera entendu qu'il s'agit de la
forme la plus forte $\mathrm{loc}(7a)$, qui implique
les deux autres. Il n'en faut pas
\deleted{\unsure{tt}} \unsure{tellement, à l'aire|tellement, à vrai dire}. Mais
\unsure{je rappelle|rappelons} que la variante \inserted{faible} \deleted{\unsure{Loc}} $(7b)$ est
\unsure{suffisante} pour faire la théorie des foncteurs
$\deleted{\unsure{f_!}} f_!$ pour les catégories $\mathrm{HOT}_W(S)$\textsuperscript{*} (cf $\overline{\underline{\mathrm{VI}}}$).
% slanted
\margin{Ici ma mémoire me joue un tour -- on n'a nullement besoin pour cela de l'\unsure{axiome} déjà très fort $\mathrm{loc}(7b)$, mais seulement de $\mathrm{Loc}(3\ \mathrm{bis})$, dont il \unsure{semble bien} qu'il soit vérifié pour tous les localiseurs fond. ``raisonnables''. \deleted{\illegible{1}}}
% archive sheet number: 311
